% Mathematical long-form supplement. Task--06: every numbered item of the original scope.
\clearpage
\section{Detailed foundations, proofs, and precise scope of the construction}\label{app:longform}
The entries below are numbered consecutively as Sections~B.1--B.89, without a second planning-number system. Each entry identifies and links to its relevant discussion in the main text. Some introductory, status, and concluding topics have no individually numbered counterpart in the main text; these are explicitly mapped to the corresponding main \emph{section}, never to an invented subsection. The main text contains the principal statements; here we spell out auxiliary arguments, hypotheses, and the precise limits of what is proved. The separate formal correspondence lists Lean declarations and their exact conditions.

\subsection{Objects and dependency of the constructions}\label{supp:1-1}
\par\noindent{\small\emph{Related main-text discussion:} \hyperref[sec:intro]{Section~\ref*{sec:intro}: Introduction} (contextual discussion; no separate main-text subsection).}\par\smallskip

The direction of dependence matters. The unary word carrier is defined before the general finite index alphabet $\Fin(b)$; the integer normal forms are built on the unary operations; positive-denominator fractions require that integer ring; and admissible three-word configurations represent those fractions once an alphabet has two designated letters. Finite child-coordinate modules are a distinct construction over the resulting rational field. Neither the space of infinite rays nor the modulated Cauchy quotient is the limit of those finite coordinate spaces by definition.

\subsection{Standard mathematical background versus proved project claims}\label{supp:1-2}
\par\noindent{\small\emph{Related main-text discussion:} \hyperref[sec:intro]{Section~\ref*{sec:intro}: Introduction} (contextual discussion; no separate main-text subsection).}\par\smallskip

The familiar settings of rooted-tree combinatorics, signed-pair completions, fraction fields, finite free modules, sequence spaces, and constructive Cauchy reals guide the terminology. A standard presentation in a textbook does not prove the project's particular definitions or its quantifiers. We use those established subjects only for context; every substantive claim about this construction is stated with the hypotheses under which the corresponding formal declaration is available.

\subsection{The scope of the main theorems}\label{supp:1-3}
\par\noindent{\small\emph{Related main-text discussion:} \hyperref[sec:intro]{Section~\ref*{sec:intro}: Introduction} (contextual discussion; no separate main-text subsection).}\par\smallskip

The main algebraic comparison is the ordered-field identification of admissible triple classes with positive-denominator fractions, and separately the decomposition $F_b=L_b\oplus S_b$ with its explicit coordinate basis. At countable arity the algebraic transverse basis survives while the chosen zero-padding maps do not preserve uniform lines. Further independent results concern prefix-boundary topology, root-preserving automorphisms and a quotient of modulated rational Cauchy sequences. This is a collection of precise results over distinct carriers, not one undifferentiated limiting construction.

\subsection{Constructive logic and scope of formal verification}\label{supp:1-4}
\par\noindent{\small\emph{Related main-text discussion:} \hyperref[sec:intro]{Section~\ref*{sec:intro}: Introduction} (contextual discussion; no separate main-text subsection).}\par\smallskip

Quotient types and structural induction are part of the background logic. The project permits the reported foundational dependencies on quotient soundness and propositional extensionality, but does not assume arbitrary classical choice in its theorem closures. An instance of an ordered-field interface is supplied with concrete decision and Archimedean data; those are mathematical premises of its generic theorems. Formal proof and an accurate mathematical interpretation of the formal statement are separate requirements; conditional completeness, inverse and order theorems remain conditional.

\subsection{Alphabet, words, incidence, and the root}\label{supp:2-1}
\par\noindent{\small\emph{Related main-text discussion:} \hyperref[main:sub-01]{Section~\ref*{main:sub-01}: Words, root and parent relation}.}\par\smallskip

The carrier is the inductively generated set of lists over $A$. This is a primitive free-word construction, rather than a graph equipped with an externally prescribed rational address. In the Stern--Brocot setting, mediants of rational fractions generate a particular binary labelling; no such rational labelling or mediant operation is assumed in the present definitions (see Section~\ref{subsec:sternbrocot}). The carrier is the inductively generated set of lists over $A$. Every list is either empty or uniquely of the form $a::v$. Consequently an immediate parent equation $w=a::v$ determines both $a$ and $v$, while the empty list cannot be a child because constructors for empty and nonempty lists are disjoint. An edge is not merely an unordered pair: it records a parent $v$, a letter $a$ and the child $a::v$. Equality of child endpoints determines an edge by injectivity of the list constructor. Thus the letter indexing of children is a bijection $A\to\{w:w\text{ is an immediate child of }v\}$ for every fixed $v$. This bijection is valid even for an empty alphabet, when both sets are empty.

\subsection{Root path existence and uniqueness}\label{supp:2-2}
\par\noindent{\small\emph{Related main-text discussion:} \hyperref[main:sub-01]{Section~\ref*{main:sub-01}: Words, root and parent relation}.}\par\smallskip

Write $E(w)$ for the directed edge list from the root to $w$. The recursion $E(\epsilon)=[]$, $E(a::v)=E(v)\mathbin{+\!+}[(v,a::v)]$ yields a directed path by induction. Conversely suppose a directed path ends at $a::v$. Its last edge is determined by the endpoint and has parent $v$. Remove that last edge and invoke induction on the path to identify the remaining list with $E(v)$. This establishes uniqueness of the \emph{edge list} rather than merely uniqueness of an endpoint sequence. In particular it does not say that an arbitrary undirected walk without reducedness constraints is unique: a walk can traverse an edge and immediately return.

\subsection{Connectivity, cycles, and arbitrarily large depth}\label{supp:2-3}
\par\noindent{\small\emph{Related main-text discussion:} \hyperref[main:sub-03]{Section~\ref*{main:sub-03}: Connectedness, cycles and depth}.}\par\smallskip

For any $u,v$, follow $E(u)$ backwards to the root and $E(v)$ forwards. This gives an undirected walk from $u$ to $v$. For each edge $e=(p,a::p)$, the set $D_e$ of descendants of $a::p$ is separated from its complement by the single edge $e$: a step between the two sides cannot use any different parent--child edge, since a child can leave its ancestor chain only through its unique parent. A purported nonempty closed walk using no edge twice would cross the cut at $e$ once, which is impossible for a closed walk. This proves the source's edge-distinct undirected no-cycle predicate; unrestricted closed backtracking walks certainly exist. When $a\in A$, the repeated words $a^n$ satisfy $d(a^n)=n$ by induction. Without an inhabitant of $A$, no unbounded-depth conclusion is asserted.

\subsection{Finite arity and its enumeration}\label{supp:2-4}
\par\noindent{\small\emph{Related main-text discussion:} \hyperref[main:sub-04]{Section~\ref*{main:sub-04}: Finite and countable branching}.}\par\smallskip

After defining $\N$, put $\Fin(b)=\{i\in\N:i<b\}$. The empty arity has no child index, since $i<0$ is impossible. The successor alphabet decomposes into the old indices $i<b$ and the new index $b$. These two cases are disjoint because $i<b$ prevents equality with $b$. Induction on $b$ therefore constructs the finite enumeration: $\operatorname{enum}(0)=[]$ and $\operatorname{enum}(b+1)$ is the old enumeration included in the successor alphabet, followed by the new index. No index is repeated, and every child index occurs. Finite coordinate sums below are folds over this \emph{proved} enumeration, not a primitive assumption that the alphabet has a prescribed cardinality.

\subsection{The countable branching object}\label{supp:2-5}
\par\noindent{\small\emph{Related main-text discussion:} \hyperref[main:sub-04]{Section~\ref*{main:sub-04}: Finite and countable branching}.}\par\smallskip

The alphabet of the $\omega$-tree is $\N$ itself. At a vertex $v$ its children are the pairwise different words $n::v$, indexed by $n\in\N$. Given any finite list of candidate children, collect their first letters; structural induction over this list yields an index not appearing in it (for example a strict upper bound on the finitely many indices). That index produces a child outside the list. This establishes the precise infinitude theorem: no finite enumeration of all children exists. It is independent of a cardinality claim for the entire set of infinite rays. The same repeated-letter depth argument gives vertices of every finite depth.

\subsection{Embeddings and finite-word exhaustion}\label{supp:2-6}
\par\noindent{\small\emph{Related main-text discussion:} \hyperref[main:sub-04]{Section~\ref*{main:sub-04}: Finite and countable branching}.}\par\smallskip

An injective alphabet map $f:A\to B$ induces an injective word map $f_*$ by entrywise substitution; a proof by induction strips the equal heads and tails. It preserves depth and incidence because $f_*(a::w)=f(a)::f_*(w)$. The specific formal maps include $\Fin(b)\hookrightarrow\Fin(b+1)$ and $\Fin(b)\hookrightarrow\N$. A map for a general inequality $b\le c$ is obtained mathematically by finitely iterating the successor inclusion, not by citing a separately declared direct map. Given $w\in\N^{<\omega}$, choose $b$ larger than every entry of the finite list $w$. Replacing each entry by its certified membership in $\Fin(b)$ gives a finite-arity preimage. The argument for an edge uses the finite list of letters in its parent and its edge symbol. Notice that the bound may be positive even for the empty word; its precise numerical value is irrelevant.

\subsection{Edge folds and their concatenation identity}\label{supp:3-1}
\par\noindent{\small\emph{Related main-text discussion:} \hyperref[main:sub-05]{Section~\ref*{main:sub-05}: Finite sums on the actual root-edge list}.}\par\smallskip

For a finite edge list $[e_1,\ldots,e_r]$, the unit fold is $1+\cdots+1$ ($r$ terms), computed by the recursion $S([])=0$ and $S(e::es)=1+S(es)$. The concatenation law $S(es\mathbin{+\!+}fs)=S(es)+S(fs)$ follows by induction on $es$: the empty case is the identity $0+S(fs)=S(fs)$ and the successor case uses associativity of $\N$ addition. This lemma is needed to justify extending the root-edge list by one edge; it also shows that edge-sum is additive under concatenating directed path segments when such concatenation is defined.

\subsection{Equality with depth}\label{supp:3-2}
\par\noindent{\small\emph{Related main-text discussion:} \hyperref[main:sub-05]{Section~\ref*{main:sub-05}: Finite sums on the actual root-edge list}.}\par\smallskip

Use the recursion $E(\epsilon)=[]$ and $E(a::w)=E(w)\mathbin{+\!+}[(w,a::w)]$. Then $S(E(\epsilon))=0=d(\epsilon)$ and, by the concatenation identity, $S(E(a::w))=S(E(w))+1$. Since $d(a::w)=d(w)+1$, induction proves $S(E(w))=d(w)$ for every $w$. This argument takes place on the actual edge list: replacing a letter preserves the sum because it does not change the number of edges, not because the vertices were independently assigned numerical coordinates.

\subsection{Ray-prefix counts are non-Cauchy}\label{supp:3-3}
\par\noindent{\small\emph{Related main-text discussion:} \hyperref[main:sub-06]{Section~\ref*{main:sub-06}: Primitive observations along rays}.}\par\smallskip

The prefix at $n$ is formed by a recursion which prepends the $n$th symbol at the successor step; thus its depth is $n$ independently of the ray. The embedded rational sequence of prefix counts therefore has values $n$. Fix a dyadic tolerance $\eps_k<1$. Given any proposed tail index $N$, take $m=N$ and $n=N+1$; their absolute difference is exactly $1$, not bounded above by $\eps_k$. Since a supplied Cauchy modulus must pass this test at every $k$, no modulus exists for this raw prefix-count sequence. This is a failure of a specified observation, not an impossibility theorem for all maps from the boundary to another number carrier.

\subsection{Unary induction and unique recursion}\label{supp:4-1}
\par\noindent{\small\emph{Related main-text discussion:} \hyperref[main:sub-07]{Section~\ref*{main:sub-07}: Construction and recursion}.}\par\smallskip

Each unary word has one of exactly two forms, $0=[]$ or $s(n)=\ast::n$. Hence any predicate holding at $0$ and propagated from $n$ to $s(n)$ holds for every element. Given a target type $X$, an element $x_0$ and an endomap $f:X\to X$, define $h(0)=x_0$ and $h(s(n))=f(h(n))$. If $g$ satisfies the same equations, the induction hypothesis $g(n)=h(n)$ implies $g(s(n))=f(g(n))=f(h(n))=h(s(n))$; uniqueness follows. This is recursion on the unary word representation, not a claim that the host logic's inductive natural-number type has been reconstructed.

\subsection{Semiring identities}\label{supp:4-2}
\par\noindent{\small\emph{Related main-text discussion:} \hyperref[main:sub-07]{Section~\ref*{main:sub-07}: Construction and recursion}.}\par\smallskip

Addition is unary concatenation. Its associativity is list-concatenation associativity, while commutativity follows from the fact that concatenating $m$ and $n$ one-letter lists yields $m+n$ copies of the unique letter in either order. Multiplication is repeated concatenation; the recursion $m\cdot0=0$, $m\cdot s(n)=m\cdot n+m$ proves $m(n+p)=mn+mp$ by induction on $p$. The identity $s(m)n=mn+n$ follows by induction on $n$. Combining the two recursions proves $mn=nm$ by induction and then right distributivity. For associativity compare $(mn)s(p)=(mn)p+mn$ with $m(ns(p))=m(np+n)=m(np)+mn$ and use induction on $p$. The zero and unit equations come from the empty and one-letter cases. No unmentioned commutative-ring properties are needed.

\subsection{Cancellation and total order}\label{supp:4-3}
\par\noindent{\small\emph{Related main-text discussion:} \hyperref[main:sub-08]{Section~\ref*{main:sub-08}: Order and cancellation}.}\par\smallskip

The equation $m+k=n+k$ between unary words implies $m=n$: induction on $k$ removes one common copy of the only letter, using injectivity of cons. Define $m\le n$ by $n=m+k$. Structural induction on the second word gives either $m\le n$ or $n<m$, so the relation is total. If $m\le n\le m$, write $n=m+p$ and $m=n+q$, hence $m=m+(p+q)$; cancellation gives $p+q=0$, forcing $p=q=0$ by the list constructors. Therefore the order is antisymmetric. Translating both sides by the same word preserves inequalities and reflects them by cancellation. The strict order $m<n$ has the witness $n=m+s(k)$, which makes $m<s(m)$ explicit.

\subsection{Quotient operation congruences}\label{supp:4-4}
\par\noindent{\small\emph{Related main-text discussion:} \hyperref[main:sub-09]{Section~\ref*{main:sub-09}: Equal-depth quotients}.}\par\smallskip

For arbitrary $A$, declare $w\sim_d v$ iff $d(w)=d(v)$ and write $[w]$ for the class. If $d(w)=d(w')$ and $d(v)=d(v')$, then the concatenations have equal lengths $d(w)+d(v)=d(w')+d(v')$. The repetition operator obeys $d(\operatorname{repeatW}(w,v))=d(w)d(v)$, by recursion on $w$; this is unchanged under both representatives. Thus addition and multiplication descend to classes. The map $[w]\mapsto d(w)$ preserves both operations and reflects equality. Given $a\in A$ it is onto by the witness $a^n$, which provides an inverse $n\mapsto[a^n]$ preserving order. If $A$ is empty, there is only the empty word and its single class; it cannot represent a positive count. The nonemptiness hypothesis is therefore not dispensable.

\subsection{Oriented differences and the congruence relation}\label{supp:5-1}
\par\noindent{\small\emph{Related main-text discussion:} \hyperref[main:sub-10]{Section~\ref*{main:sub-10}: Difference pairs and equivalence}.}\par\smallskip

On $\N\times\N$ set $(m,n)\sim(p,q)$ iff $m+q=p+n$. The relation is reflexive and symmetric by commutativity of addition. For transitivity, suppose also $p+s=r+q$. Add the two given equalities and reassociate: $m+q+p+s=p+n+r+q$. Cancellation of the common $p+q$ yields $m+s=r+n$. The equivalence relation encodes a formal difference $m-n$ but uses no primitive negative edge weights and no preconstructed signed number type.

\subsection{Normalization and exactness}\label{supp:5-2}
\par\noindent{\small\emph{Related main-text discussion:} \hyperref[main:sub-11]{Section~\ref*{main:sub-11}: Normal forms}.}\par\smallskip

Simultaneously remove a successor from both components of $(m,n)$ until one becomes zero. This procedure terminates by recursion on unary words and produces $(m',n')$ with $m'=0$ or $n'=0$. Each removal preserves the difference relation: $(s(m),s(n))\sim(m,n)$ holds because $s(m)+n=m+s(n)=s(m+n)$. Iterating this step leaves the represented formal difference unchanged. Two normalized pairs representing the same difference cannot have opposite nonzero signs: if $(a,0)\sim(0,b)$, then $a+b=0$, forcing both to vanish. If they have the same sign, unary cancellation gives equality of the nonzero components. Hence normalization is exact and unique, and the type $\Z$ of normalized pairs is a concrete carrier in bijection with formal difference classes. This is \emph{not} literally implemented as a quotient type for $\Z$.

\subsection{Integer operations and zero divisors}\label{supp:5-3}
\par\noindent{\small\emph{Related main-text discussion:} \hyperref[main:sub-12]{Section~\ref*{main:sub-12}: Arithmetic and integral-domain laws}.}\par\smallskip

On formal pairs take $(m,n)+(p,q)=(m+p,n+q)$, $-(m,n)=(n,m)$ and $(m,n)(p,q)=(mp+nq,mq+np)$, then normalize. For product congruence in the first argument, multiply $m+q=p+n$ by the appropriate two components of the second difference and use distributivity to rewrite the same cross-product equation; similarly in the other argument. Thus all operations descend to normal forms. Associativity, commutativity and distributivity follow by expanding the corresponding polynomial identities over $\N$ and using normalization exactness. For the zero-product law: if the normalized product vanishes, then the represented formal product difference is zero; inspect the two normalized signs. If neither factor is zero, their nonzero components have products strictly positive by unary product positivity, contradicting a zero difference. This also yields cancellation of a nonzero factor by applying zero-product to $x(y-z)=0$.

\subsection{Integer order and natural embedding}\label{supp:5-4}
\par\noindent{\small\emph{Related main-text discussion:} \hyperref[main:sub-13]{Section~\ref*{main:sub-13}: Natural embedding and order}.}\par\smallskip

A normalized integer is nonnegative exactly when its negative component is zero. Thus $x\le y$ means that the normalized representative of $y-x$ has zero negative component. Sign normalization makes this decidable and total; transitivity follows from closure of nonnegative elements under addition. The map $n\mapsto(n,0)$ preserves zero and one and satisfies $(m,0)+(n,0)=(m+n,0)$ and $(m,0)(n,0)=(mn,0)$, hence is a semiring homomorphism. If $(m,0)=(n,0)$, their positive coordinates agree, proving injectivity. The induced order agrees with the unary prefix/addition order, because $(n,0)-(m,0)$ is nonnegative exactly when $n=m+k$ for some $k$.

\subsection{Admissibility and positive denominators}\label{supp:6-1}
\par\noindent{\small\emph{Related main-text discussion:} \hyperref[main:sub-14]{Section~\ref*{main:sub-14}: Admissible data}.}\par\smallskip

A triple $t=(P_1,P_2,P_3)$ is admissible when its words are pairwise different and the positive difference $d(P_2)-d(P_3)>0$ is supplied. Its signed numerator is $a(t)=d(P_1)-d(P_2)$ in the normalized integer carrier and its denominator is $D(t)=d(P_2)-d(P_3)$ in $\N$. Pairwise inequality of words is stronger than a numerical inequality between lengths when the alphabet contains two letters: different letters can distinguish paths of equal length. This distinction is precisely why branching supplies witnesses even for a zero numerator. The denominator, by contrast, is numerical and is strictly positive by the defining hypothesis; every rational operation must preserve that condition.

\subsection{Fraction equivalence and transitivity}\label{supp:6-2}
\par\noindent{\small\emph{Related main-text discussion:} \hyperref[main:sub-15]{Section~\ref*{main:sub-15}: The auxiliary fraction relation}.}\par\smallskip

Raw fractions are pairs $(a,D)$ with $a\in\Z$, $D\in(\N)_{>0}$. Define $(a,D)\sim(c,E)$ by $a\overline E=c\overline D$, with bars denoting the natural-to-integer embedding. Reflexivity and symmetry are immediate. If $a\overline E=c\overline D$ and $c\overline F=g\overline E$, multiply the first equation by $\overline F$ and the second by $\overline D$. The results are $a\overline E\overline F=c\overline D\overline F=g\overline E\overline D$. Since $\overline E\ne0$, integral-domain cancellation yields $a\overline F=g\overline D$. Notice that transitivity \emph{uses} positivity of the middle denominator; cross multiplication with an arbitrary zero denominator would not define the intended equivalence relation.

\subsection{Every signed fraction has a two-letter triple witness}\label{supp:6-3}
\par\noindent{\small\emph{Related main-text discussion:} \hyperref[main:sub-17]{Section~\ref*{main:sub-17}: Universal triple witnesses}.}\par\smallskip

Choose distinct $\alpha_0,\alpha_1$. Express a signed numerator as the normalized difference $p-q$, and let $D>0$. The three words
\[
 P_1=\alpha_0^{p+2+D},\qquad
 P_2=\alpha_1^{q+2+D},\qquad
 P_3=\alpha_0::\alpha_1^{q+1}
\]
have depths $p+2+D$, $q+2+D$ and $q+2$. Their numerator is $p-q$ and denominator is $D$. The first letters distinguish $P_1$ from $P_2$, and $P_2$ from $P_3$. The first two letters distinguish $P_1$ from $P_3$, even if their depths happen to be the same; all prescribed exponents ensure these letters exist. This last check is essential when the numerator is zero or negative. Hence every positive-denominator fraction has a representative in the triple type with two letters. When the alphabet is unary, unequal words have unequal lengths: admissibility then excludes both numerator $0$ and numerator $-D$ (value $-1$). The exact complement of the unary image is not claimed as a kernel theorem.

\subsection{Congruence and the transported field operations}\label{supp:6-4}
\par\noindent{\small\emph{Related main-text discussion:} \hyperref[main:sub-18]{Section~\ref*{main:sub-18}: Transported field structure and order}.}\par\smallskip

For $(a,D)\sim(a',D')$ and $(c,E)\sim(c',E')$, multiplication respects the relation since
\[(ac)\overline{D'E'}=(a\overline{D'})(c\overline{E'})=(a'\overline D)(c'\overline E)=(a'c')\overline{DE}.\]
Addition respects it by two successive one-variable replacements: using $a\overline{D'}=a'\overline D$,
\[(a\overline E+c\overline D)\overline{D'E}=(a'\overline E+c\overline{D'})\overline{DE},\]
and similarly for replacing $(c,E)$ by $(c',E')$. Denominators remain positive under multiplication. The usual three-fraction common-denominator expansion proves associativity and distributivity; $a/D+(-a)/D=0$ because the numerator is zero, and the units are $0/1,1/1$. A nonzero rational fraction has numerator of definite sign in $\Z$; the sign-normalized reciprocal has positive denominator, is independent of representatives and multiplies back to $1$. The \emph{path-triple} field operations are then \emph{transported} through the bijection with this auxiliary fraction field. No independent incidence-defined field multiplication is postulated.

\subsection{The rational order and explicit Archimedean data}\label{supp:6-5}
\par\noindent{\small\emph{Related main-text discussion:} \hyperref[main:sub-18]{Section~\ref*{main:sub-18}: Transported field structure and order}.}\par\smallskip

Positive denominators permit cross-product comparison:
\[(a,D)\le(c,E)\iff a\overline E\le c\overline D.\]
To check representative invariance, compare $c\overline D'-a'\overline E$ and $c\overline D-a\overline E$: multiplication by the positive denominators $D$ and $D'$ converts one to the other. Positivity and integral-domain cancellation preserve the sign. The numerator-normal-form sign decision proves that any two fractions are comparable; multiplication by a nonnegative fraction respects order by multiplying the cross-product inequality by nonnegative integers. For a constructive Archimedean upper bound, write $a=p-q$ and use $a\le p$, $D\ge1$, obtaining $a/D\le p$ when the comparisons are performed with positive denominators. This gives an explicit bounding count. Dyadic fractions $2^{-k}$ are well defined since $2^k>0$, their successive halves sum to the preceding precision, and Archimedean comparison yields a precision eventually below any given positive rational $\delta$.

\subsection{Exact inverse maps to the auxiliary fraction field}\label{supp:6-6}
\par\noindent{\small\emph{Related main-text discussion:} \hyperref[main:sub-19]{Section~\ref*{main:sub-19}: Identification with the internal fraction field}.}\par\smallskip

Let $\Phi([t])=[(a(t),D(t))]$. Quotient congruence makes $\Phi$ independent of the triple representative. For a positive-denominator raw fraction $f$, use the preceding two-letter witness $t_f$ and put $\Psi([f])=[t_f]$. If $f\sim g$, their witnesses are equivalent by the same cross-product equation; hence $\Psi$ is independent of the chosen raw fraction. By construction $\Phi\Psi([f])=[f]$. Conversely $\Psi\Phi([t])$ has the same numerator and denominator fraction as $t$; the equivalence defining the triple quotient gives $\Psi\Phi([t])=[t]$. Both maps respect all operations because these operations were defined by transport. Order reflection and preservation follow from the common positive-denominator formula, yielding an \emph{ordered} field isomorphism rather than merely a bijection.

\subsection{The precise contribution of branching}\label{supp:6-7}
\par\noindent{\small\emph{Related main-text discussion:} \hyperref[main:sub-17]{Section~\ref*{main:sub-17}: Universal triple witnesses}.}\par\smallskip

Any two words of equal depth have the same path sum. Therefore the rational value attached to a triple depends only on the three lengths, even though the admissibility condition remembers inequality of the words themselves. Distinct letters are useful for producing three different words with otherwise prescribed lengths; they are not necessary for defining formal differences or for algebraic arithmetic in the auxiliary fraction field. In particular the unary failure is a restriction of the \emph{particular triple representation}, not a mathematical impossibility of a rational field based on unary natural numbers.

\subsection{Arity transfer as composition through fractions}\label{supp:7-1}
\par\noindent{\small\emph{Related main-text discussion:} \hyperref[main:sub-20]{Section~\ref*{main:sub-20}: Comparison across alphabets}.}\par\smallskip

For admissible alphabets $A,B$, define $T_{A,B}=\Psi_B\Phi_A$. Since $\Phi_A$ and $\Psi_B$ are inverse to the respective fraction-field comparisons, $T_{A,B}$ is a field isomorphism. The ordered version uses the separately proved preservation of the transported order in both directions. This construction makes no use of a bijection $A\to B$; indeed such a bijection need not exist for different finite arities. It instead compares the numerical \emph{quotients} of the two distinct underlying word systems.

\subsection{Identity, composition and embeddings}\label{supp:7-2}
\par\noindent{\small\emph{Related main-text discussion:} \hyperref[main:sub-20]{Section~\ref*{main:sub-20}: Comparison across alphabets}.}\par\smallskip

For admissible $A,B,C$ compute
\[
T_{B,C}T_{A,B}=(\Psi_C\Phi_B)(\Psi_B\Phi_A)=\Psi_C(\Phi_B\Psi_B)\Phi_A=\Psi_C\Phi_A=T_{A,C}.
\]
Likewise $T_{A,A}=\Psi_A\Phi_A=\mathrm{id}$. These equalities are identities of field maps, not of individual words. If $n\in\N$ embeds as $n/1$ in the fraction carrier, then $\Phi_A$ sends its path-triple presentation to $n/1$ and $\Psi_B$ sends that fraction to its chosen witness in $B$; compatibility with the natural embedding follows. Explicit letter-choice independence at the quotient level must not be misread as equality of two different triple witnesses prior to quotienting.

\subsection{Arithmetic versus local coordinate rank}\label{supp:7-3}
\par\noindent{\small\emph{Related main-text discussion:} \hyperref[main:sub-21]{Section~\ref*{main:sub-21}: What the comparison does not identify}.}\par\smallskip

An isomorphism $\Q(A)\simeq\Q(B)$ concerns the abstract field of represented rational values; no vector-space dimension is attached to an abstract field by this statement. In Section~\ref{sec:finite}, $b$ instead indexes a separate free coordinate module with $b$ coordinates. Its rank varies with $b$ even when all coefficient fields are transported to a common one. The two assertions are therefore compatible; they are not competing dimension claims.

\subsection{A finite free coordinate module}\label{supp:8-1}
\par\noindent{\small\emph{Related main-text discussion:} \hyperref[main:sub-22]{Section~\ref*{main:sub-22}: Coordinates and the chosen bilinear form}.}\par\smallskip

For positive $b$ let $F_b$ consist of functions $\Fin(b)\to K$. All operations are coordinatewise. Evaluation at a child index is linear because $(x+y)_i=x_i+y_i$ and $(ax)_i=ax_i$. Let $e_i(j)$ equal $1$ when $i=j$ and $0$ otherwise. Then $x=\sum_{i<b}x_ie_i$ by evaluating both sides at an arbitrary $j$ and observing that only $i=j$ contributes. Conversely, if $\sum_i a_ie_i=0$, evaluating at $j$ gives $a_j=0$, proving the familiar coordinate family a basis. These are finite sums over the internally enumerated alphabet, not infinite series.

\subsection{Finite coordinate sums and symmetric bilinearity}\label{supp:8-2}
\par\noindent{\small\emph{Related main-text discussion:} \hyperref[main:sub-22]{Section~\ref*{main:sub-22}: Coordinates and the chosen bilinear form}.}\par\smallskip

Write $\sigma(x)=\sum_i x_i$ and $B(x,y)=\sum_i x_iy_i$. Distributivity in the coefficient field gives $B(x+x',y)=\sum_i(x_i+x'_i)y_i=B(x,y)+B(x',y)$. Likewise $B(ax,y)=aB(x,y)$ because finite sums permit factoring a constant scalar. Commutativity of field multiplication gives $B(x,y)=B(y,x)$, so additivity and homogeneity in the second argument follow by symmetry. This proves all four Task--06 bilinearity identities, including their independent right-sided versions. A symmetric bilinear form over an arbitrary abstract field is not asserted to be a metric or a positive-definite inner product.

\subsection{The uniform line and the zero-sum space}\label{supp:8-3}
\par\noindent{\small\emph{Related main-text discussion:} \hyperref[main:sub-23]{Section~\ref*{main:sub-23}: Longitudinal and transverse subspaces}.}\par\smallskip

Put $u=(1,\ldots,1)$, $L_b=\{au:a\in K\}$, and $S_b=\ker\sigma$. Closure of $L_b$ under addition and scalar multiplication follows from $(a+a')u=au+a'u$ and $c(au)=(ca)u$. Closure of $S_b$ follows because $\sigma(x+y)=\sigma(x)+\sigma(y)$ and $\sigma(ax)=a\sigma(x)$. Nonemptiness of the child alphabet when $b>0$ gives a distinguished coordinate at which $u$ has entry $1$, hence $u\ne0$. For $b=0$, the uniform function on the empty type is the zero vector, and the rank-one assertion would be false; this is why the positive-arity qualification is essential.

\subsection{Algebraic orthogonality}\label{supp:8-4}
\par\noindent{\small\emph{Related main-text discussion:} \hyperref[main:sub-26]{Section~\ref*{main:sub-26}: Orthogonality and the splitting of the bilinear form}.}\par\smallskip

For $s\in S_b$, the bilinear form with the uniform vector satisfies $B(u,s)=\sum_i 1\cdot s_i=\sigma(s)=0$. If $l=au\in L_b$, bilinearity gives $B(l,s)=aB(u,s)=0$. Conversely, no assertion is needed that the difference basis vectors $e_j-e_0$ are pairwise orthogonal: indeed $B(e_j-e_0,e_k-e_0)=1$ for distinct positive indices $j,k$ in characteristic zero, so they are usually \emph{not} mutually orthogonal. The proved orthogonality is that of the two \emph{subspaces} under this particular bilinear form.

\subsection{A basis for the uniform line}\label{supp:8-5}
\par\noindent{\small\emph{Related main-text discussion:} \hyperref[main:sub-23]{Section~\ref*{main:sub-23}: Longitudinal and transverse subspaces}.}\par\smallskip

The singleton family $\{u\}$ spans $L_b$ by definition. If $au=0$, evaluation at any index (which exists for $b>0$) yields $a=0$; hence the family is linearly independent. This is the exact rank-one basis claim. In the concrete instance $K=\Q(A)$, the required no-collapse property follows from the constructed rational field and its characteristic zero; no completion is used to obtain the basis.

\subsection{The transverse basis with $b-1$ members}\label{supp:8-6}
\par\noindent{\small\emph{Related main-text discussion:} \hyperref[main:sub-23]{Section~\ref*{main:sub-23}: Longitudinal and transverse subspaces}.}\par\smallskip

For $b\ge1$, index positive coordinates by $j=1,\ldots,b-1$ and set $v_j=e_j-e_0$. Their sums vanish. In any finite dependence $\sum_j a_jv_j=0$, the $j$th positive coordinate equals $a_j$, forcing every coefficient to be zero. For $x\in S_b$ the equality $\sigma(x)=0$ forces $x_0=-\sum_{j>0}x_j$, giving the finite expansion $x=\sum_{j>0}x_jv_j$. The positive indices form the successor-predecessor child alphabet in the source, so the basis has exactly $b-1$ elements. This constructs the basis, rather than inferring a dimension from a pre-existing rank--nullity theorem.

\subsection{Complementary projections and the direct sum}\label{supp:8-7}
\par\noindent{\small\emph{Related main-text discussion:} \hyperref[main:sub-25]{Section~\ref*{main:sub-25}: Complementary linear projections and their images}.}\par\smallskip

Assume $b>0$ and that its numeral $b1_K$ is invertible (in particular under the project's characteristic-zero field interface). Define $P_Lx=(b1_K)^{-1}\sigma(x)u$ and $P_Sx=x-P_Lx$. Then $\sigma(P_Sx)=\sigma(x)-b(b1_K)^{-1}\sigma(x)=0$, while $P_Lx\in L_b$. Hence $x=P_Lx+P_Sx$. If $au\in L_b\cap S_b$, then $ba=0$ and invertibility gives $a=0$. This proves uniqueness. Linearity is obtained by applying linearity of $\sigma$ inside $P_L$ and subtracting for $P_S$. Since $\sigma(P_Lx)=\sigma(x)$, $P_L^2=P_L$; and since $P_Sx\in\ker\sigma$, $P_LP_S=0$. The identity $P_S=I-P_L$ then yields $P_S^2=P_S$ and $P_SP_L=0$. If $x\in L_b$, then $P_Lx=x$ and $P_Sx=0$; if $x\in S_b$, then $P_Lx=0$ and $P_Sx=x$. These identities prove $\operatorname{im}P_L=L_b=\ker P_S$ and $\operatorname{im}P_S=S_b=\ker P_L$. The bare equality $x=P_Lx+P_Sx$ is definitional, whereas membership in the stated subspaces and their directness use the hypotheses.

\subsection{Unary arity and permutation symmetry}\label{supp:8-8}
\par\noindent{\small\emph{Related main-text discussion:} \hyperref[main:sub-27]{Section~\ref*{main:sub-27}: Permutations and the unary exception}.}\par\smallskip

When $b=1$, $\sigma(x)=x_0$, so $S_1=\{0\}$, $L_1=F_1$ and $P_L=I$, $P_S=0$. More generally any permutation $\pi$ of the finite child alphabet defines a reindexing $(\pi x)_i=x_{\pi^{-1}i}$. Reindexing a finite enumerated sum leaves its value unchanged, so $\sigma(\pi x)=\sigma(x)$ and $B(\pi x,\pi y)=B(x,y)$. Since $\pi u=u$, both $L_b$ and $S_b$ are preserved; the projections also commute with the permutation action. These are algebraic symmetries of child coordinates, not statements about all tree automorphisms.

\subsection{What $b-1$ measures}\label{supp:8-9}
\par\noindent{\small\emph{Related main-text discussion:} \hyperref[main:sub-23]{Section~\ref*{main:sub-23}: Longitudinal and transverse subspaces}.}\par\smallskip

The integer $b-1$ is the cardinality of an explicitly constructed finite basis of the kernel of a coordinate-sum linear map. This establishes an algebraic free rank over the chosen coefficient field. It does not define, or compute, the topological covering dimension of the finite tree, its ray boundary, a differentiable manifold or a physical model. Any such interpretation would require further structures and new theorems absent from these hypotheses.

\subsection{The algebraic direct sum and finite coordinate sums}\label{supp:9-1}
\par\noindent{\small\emph{Related main-text discussion:} \hyperref[main:sub-28]{Section~\ref*{main:sub-28}: Finite support and coordinate sums}.}\par\smallskip

Let $K^{(\N)}$ denote functions $x:\N\to K$ with finite support. Such a function admits a finite index bound $N_x$ beyond which all coordinates vanish; choose a common upper bound for finitely many such functions when adding or taking finite linear combinations. Define $\Sigma x$ as the sum of $x_i$ for $i\le N_x$. Independence of the bound follows because every additional summand is zero. Thus $\Sigma(x+y)=\Sigma x+\Sigma y$ and $\Sigma(ax)=a\Sigma x$. The transverse subspace $S_\omega=\ker\Sigma$ is therefore an algebraic submodule. No topology, infinite sum or limit of partial sums is implicit in this definition.

\subsection{A countably indexed algebraic basis}\label{supp:9-2}
\par\noindent{\small\emph{Related main-text discussion:} \hyperref[main:sub-28]{Section~\ref*{main:sub-28}: Finite support and coordinate sums}.}\par\smallskip

Set $v_i=e_{i+1}-e_0$ for every $i\in\N$. A finite relation $\sum_{i\in I}a_iv_i=0$ gives $a_i=0$ at coordinate $i+1$ for each $i\in I$, proving independence. For $x\in S_\omega$ with support within $\{0,\ldots,N\}$, the finite sum $\sum_{i<N}x_{i+1}v_i$ reproduces all positive coordinates and has zeroth coordinate $-\sum_{j=1}^{N}x_j=x_0$. Hence it equals $x$. This is algebraic spanning by finite combinations; replacing it with an infinite series and a convergence claim would be unjustified.

\subsection{Padding embeddings and compatibility}\label{supp:9-3}
\par\noindent{\small\emph{Related main-text discussion:} \hyperref[main:sub-29]{Section~\ref*{main:sub-29}: Compatible embeddings}.}\par\smallskip

The zero-padding map $\iota_b:F_b\to F_{b+1}$ retains coordinates $i<b$ and sets the new coordinate $b$ to zero. Coordinatewise evaluation gives injectivity, and the same formulas show additivity and scalar homogeneity. Since appending a zero does not change a finite sum, $\Sigma_{b+1}(\iota_bx)=\Sigma_b(x)$, so $\iota_b(S_b)\subseteq S_{b+1}$. Its iterates into $K^{(\N)}$ preserve the same zero-sum relation. The basis vector $e_j-e_0$ is sent to the same named coordinate difference, proving compatibility on the transverse basis family.

\subsection{The zero-padding obstruction for uniform lines}\label{supp:9-4}
\par\noindent{\small\emph{Related main-text discussion:} \hyperref[main:sub-30]{Section~\ref*{main:sub-30}: Exact limiting conclusion}.}\par\smallskip

For $b>0$, $\iota_b(u_b)$ has coordinate $0$ equal to $1$ but coordinate $b$ equal to $0$. If $\iota_b(u_b)=au_{b+1}$, evaluation first at $0$ and then at $b$ yields $a=1$ and $a=0$, contradicting the field's nontriviality. Likewise the everywhere-one function on $\N$ is not finitely supported: for any proposed finite bound $N$, its value at $N+1$ is still $1\ne0$. Therefore no element of the finite-support carrier equals that function. These are obstructions relative to \emph{zero padding and finite support}; they do not address enlarged ambient vector spaces or different embeddings.

\subsection{Precise statement of the countable result}\label{supp:9-5}
\par\noindent{\small\emph{Related main-text discussion:} \hyperref[main:sub-30]{Section~\ref*{main:sub-30}: Exact limiting conclusion}.}\par\smallskip

The compatible transverse injections and their finite-basis vectors describe the algebraic directed union within $S_\omega$; the full latter space is reached because every element has finite support. The uniform lines do not define a parallel compatible direct system under these inclusions, so the finite decompositions $F_b=L_b\oplus S_b$ do not themselves assemble into a corresponding longitudinal-plus-transverse decomposition of $K^{(\N)}$. There is no defined ``limit line'' here to which nonexistence could be attributed without first specifying an ambient space and its comparison maps.

\subsection{Rays, coherent prefixes and infinite depth}\label{supp:10-1}
\par\noindent{\small\emph{Related main-text discussion:} \hyperref[main:sub-31]{Section~\ref*{main:sub-31}: Rays and finite prefixes}.}\par\smallskip

A ray is a function $\xi:\N\to A$; its prefix $\xi|_n$ is the finite word constructed by the recursion $\xi|_0=\epsilon$ and $\xi|_{n+1}=\xi(n)::(\xi|_n)$. Each successor is a parent--child step, so the directed root path to $\xi|_n$ is exactly the first $n$ edges of the ray. Induction gives $d(\xi|_n)=n$. If $\xi|_n=\eta|_n$ for every $n$, then at $n+1$ equality of the heads gives $\xi(n)=\eta(n)$; functional extensionality yields $\xi=\eta$. The prefixes have arbitrarily large depth for any existing ray. Existence of a ray is not claimed for the empty alphabet, where no function from the nonempty unary index type into the alphabet exists.

\subsection{Finite-arity ray inclusions are not exhaustive}\label{supp:10-2}
\par\noindent{\small\emph{Related main-text discussion:} \hyperref[main:sub-32]{Section~\ref*{main:sub-32}: Noncommutation with finite-arity unions}.}\par\smallskip

A ray $r:\N\to\Fin(b)$ yields a ray in the $\omega$-tree by composing with the index inclusion. Equality of finite prefixes is preserved because inclusion acts symbolwise. To show this does not exhaust the countable boundary, take $\iota(n)=n$. If $\iota$ came from a fixed $r:\N\to\Fin(b)$, then evaluating at $n=b$ would give $b=(r(b)).\mathrm{idx}<b$, a contradiction. Unlike the finite-word exhaustion proof, one cannot choose a common finite upper bound for an \emph{infinite} list of ray symbols. This is the precise noncommutation between the finite-arity union and passing to infinite rays.

\subsection{Cylinders and their clopen structure}\label{supp:10-3}
\par\noindent{\small\emph{Related main-text discussion:} \hyperref[main:sub-33]{Section~\ref*{main:sub-33}: Cylinders and topology}.}\par\smallskip

For $w\in A^{<\omega}$ let $C(w)=\{\xi:w\preceq\xi\}$, where the prefix relation uses the source's orientation. Every $\xi\in C(w)$ has its length-$d(w)$ prefix fixed, and the cylinder of that prefix is contained in $C(w)$; hence $C(w)$ is open in the prefix topology. If $\xi\notin C(w)$, any ray $\eta$ sharing $\xi|_{d(w)}$ is also outside $C(w)$, because otherwise equality of their initial $d(w)$-blocks would force $\xi\in C(w)$. Thus $\rays_A\setminus C(w)$ is open. This proof does not decide equality of arbitrary letters; it uses only the consequences of having a given prefix and the negation of the cylinder membership. Intersections of cylinders containing the same ray contain a sufficiently long common-prefix cylinder, so the cylinders form a basis.

\subsection{Finite witness separation and constructive equality}\label{supp:10-4}
\par\noindent{\small\emph{Related main-text discussion:} \hyperref[main:sub-33]{Section~\ref*{main:sub-33}: Cylinders and topology}.}\par\smallskip

If $\xi(n)\ne\eta(n)$ is supplied, their prefixes of length $n+1$ are incompatible: a hypothetical ray in both corresponding cylinders would have its $n$th letter equal to both differing letters. These open cylinders separate $\xi$ and $\eta$, and each is also closed. From merely $\neg(\xi=\eta)$, however, extracting a particular disagreement index is not constructively automatic; the source only establishes a double-negated existence of such an index under its hypotheses. The resulting separation theorem is therefore expressed with an explicit disagreement witness. Classical assertions about product Cantor or Baire spaces are not imported as proofs of the source-specific prefix topology.

\subsection{Connected sets, disconnectedness and isolation}\label{supp:10-5}
\par\noindent{\small\emph{Related main-text discussion:} \hyperref[main:sub-34]{Section~\ref*{main:sub-34}: Connectedness and isolated points}.}\par\smallskip

Assume decidable equality on the alphabet $A$, as in the source's \texttt{Ray.}\allowbreak\texttt{connected\_}\allowbreak\texttt{subsingleton} declaration with the explicit hypothesis \texttt{[DecidableEq $A$]}. Under its connectedness predicate, if a subset contains two rays with a \emph{witnessed} differing coordinate, the complementary clopen cylinder pieces separate the subset. For arbitrary $\xi,\eta$ in a connected subset, decide $\xi(n)=\eta(n)$ for each $n$. A negative answer supplies the forbidden separation; therefore equality holds for each $n$, and function extensionality yields $\xi=\eta$. This step never infers a disagreement index solely from the negation of equality of infinite rays. For an alphabet with two distinguished different letters, fix any $\xi$ in a cylinder determined by a finite prefix; replace one symbol at an index after that prefix by the other designated letter. The resulting ray is different but remains in the same cylinder. Hence no basic neighborhood isolates $\xi$, and consequently no prefix-open neighborhood does. In the unary case all rays agree coordinatewise, producing the exceptional singleton boundary.

\subsection{Countable cylinder noncompactness and nonlocal compactness}\label{supp:10-6}
\par\noindent{\small\emph{Related main-text discussion:} \hyperref[main:sub-35]{Section~\ref*{main:sub-35}: Countable-arity compactness obstruction}.}\par\smallskip

Fix $w$ and cover $C(w)$ by the family $\{C(a::w):a\in\N\}$. Every ray in $C(w)$ has some next letter, hence lies in one member. A finite subfamily uses only finitely many such indices; choose a larger index $b$ and a ray beginning with $b::w$. That ray belongs to $C(w)$ but avoids the chosen subfamily. This refutes the source's finite-subcover compactness. If a point admitted a compact neighborhood $H$, the basis property would give a cylinder $C(w)$ contained in $H$. Since $C(w)$ is closed, an open cover of it extends to one of $H$ by adding its open complement. Compactness of $H$ would yield a finite subcover of the cylinder, contradiction. The argument concerns countable branching; it gives no formal compactness certificate for an arbitrary finite-alphabet boundary.

\subsection{Precisely conditional chart obstructions}\label{supp:10-7}
\par\noindent{\small\emph{Related main-text discussion:} \hyperref[main:sub-36]{Section~\ref*{main:sub-36}: Manifold-chart questions}.}\par\smallskip

A chart in the source is a homeomorphism of specified open subspaces satisfying the given continuity conditions. If its target has the property that every nonempty relevant open set contains a connected subset with two different points, the inverse chart would carry those two points to a connected subset of the boundary, contradicting connected-subsingleton. In the target $\mathrm{Unit}$, a nonunary ray boundary has no chart whose domain neighborhood must be mapped bijectively to a singleton, because that would isolate its sole preimage; the unary one-point boundary admits the trivial chart. For the constructed Cauchy real-line topology, the hypothesis producing connected target pairs is \texttt{IntervalsConnected $K$}. The resulting no-chart theorem is conditional on that hypothesis. No unproved connectedness of intervals or all positive-dimensional Euclidean charts is silently assumed.

\subsection{The limits of numerical ray encodings}\label{supp:10-8}
\par\noindent{\small\emph{Related main-text discussion:} \hyperref[main:sub-37]{Section~\ref*{main:sub-37}: Optional numerical encodings}.}\par\smallskip

A letter-to-digit assignment followed by an infinite dyadic expansion defines an \emph{additional} map out of a boundary, not its native topology. In the particular formal binary encoding two distinct rays can represent the same numerical class; this is a direct counterexample to injectivity of that map. Choosing different distinguished letters can change the encoded value. Separately, if an encoding were invariant under the transitive levelwise ray action of Section~\ref{sec:automorphisms}, then the invariance theorem would force it to be constant. These statements neither prohibit noninvariant encodings nor construct any canonical homeomorphism between the prefix boundary and the numerical completion.

\subsection{Levelwise permutations and induced automorphisms}\label{supp:11-1}
\par\noindent{\small\emph{Related main-text discussion:} \hyperref[main:sub-38]{Section~\ref*{main:sub-38}: Levelwise actions}.}\par\smallskip

Choose for each depth $n$ a permutation $\pi_n:A\to A$ with an explicitly given inverse. On a word, transform each letter by the permutation indexed by the depth of its suffix (equivalently by its distance from the root). The root is fixed. The identity $T(a::w)=\pi_{d(w)}(a)::T(w)$ shows that every parent--child edge is mapped to another such edge and that depth is preserved. The same recursion with $\pi_n^{-1}$ gives an inverse word map, by induction on word length. On rays, $(T\xi)(n)=\pi_n(\xi(n))$, and the finite-prefix identity is verified by induction. These automorphisms form a particular subgroup; the construction does not claim exhaustiveness.

\subsection{An explicit transitive action on rays}\label{supp:11-2}
\par\noindent{\small\emph{Related main-text discussion:} \hyperref[main:sub-38]{Section~\ref*{main:sub-38}: Levelwise actions}.}\par\smallskip

For given rays $\xi,\eta$ and decidable equality on $A$, at each $n$ define the involution swapping $\xi(n)$ and $\eta(n)$ and fixing all other letters. Its being a permutation follows by checking the two swapped cases and the fixed case. The induced action sends the $n$th coordinate of $\xi$ to $\eta(n)$ for every $n$, hence sends the whole ray to $\eta$. The argument constructs a \emph{level-dependent} permutation family; it does not need the same permutation at every level or a global classification of rooted-tree automorphisms.

\subsection{Invariant functions are constant}\label{supp:11-3}
\par\noindent{\small\emph{Related main-text discussion:} \hyperref[main:sub-38]{Section~\ref*{main:sub-38}: Levelwise actions}.}\par\smallskip

Let $f:\rays_A\to Y$ obey $f(T\xi)=f(\xi)$ for every member $T$ of the constructed levelwise action. Given $\xi,\eta$, use the explicit swapping family from the preceding item. Then $f(\eta)=f(T\xi)=f(\xi)$. This proves constancy without any order, topology or algebraic structure on $Y$. Conversely the theorem says nothing about maps that are not invariant under those permutations; in particular it does not rule out label-dependent numerical encodings.

\subsection{A subtree-supported automorphism is not levelwise}\label{supp:11-4}
\par\noindent{\small\emph{Related main-text discussion:} \hyperref[main:sub-39]{Section~\ref*{main:sub-39}: A non-levelwise automorphism}.}\par\smallskip

Choose two different child letters and one distinguished parent address. Swap the descendant branches beginning with the chosen two letters \emph{below that address}, while leaving the analogous branches below another address at the same depth unchanged. The swap is an involution, so it is bijective; each altered word keeps its ancestor chain up to the distinguished address and has the same tail structure thereafter. Hence adjacency and root are preserved. Suppose it were induced by a single permutation $\pi_n$ at the affected depth. The behavior below the distinguished address would force $\pi_n$ to swap the selected letters, whereas the fixed behavior below the second address would force $\pi_n$ to fix them, contradiction. This refutes the levelwise-completeness assertion; it does not supply a classification of all automorphisms.

\subsection{The base interface and dyadic precision}\label{supp:12-1}
\par\noindent{\small\emph{Related main-text discussion:} \hyperref[main:sub-40]{Section~\ref*{main:sub-40}: Precision and modulated sequences}.}\par\smallskip

The base $K$ is an algebraic characteristic-zero field with a nonnegative cone satisfying closure, totality and antisymmetry, a \emph{decision procedure} for that cone, and an explicitly supplied Archimedean bound $a\mapsto B(a)\in\N$. The weak order $x\le y$ means $\operatorname{Nonneg}(y-x)$ and its strict counterpart is the negation of the reversed weak order. Thus comparisons between \emph{rational} approximants are decidable without deciding general propositions about Cauchy classes. Define $\eps_k=2^{-k}$ using inverses of positive powers of two. Positivity, the halving identity $2\eps_{k+1}=\eps_k$, and antitonicity in $k$ follow from the ordered-field laws. For $\delta>0$, choose an Archimedean count exceeding $1/\delta$ and then a power of two beyond that count; its reciprocal is less than $\delta$. This justifies all later changes to smaller dyadic error budgets.

\subsection{Bare Cauchy, modulated Cauchy and regularity}\label{supp:12-2}
\par\noindent{\small\emph{Related main-text discussion:} \hyperref[main:sub-40]{Section~\ref*{main:sub-40}: Precision and modulated sequences}.}\par\smallskip

A bare Cauchy predicate supplies, for each positive $\delta$, \emph{some} tail $N(\delta)$ as a proposition. A modulated datum additionally stores one function $\mu(k)$ such that $m,n\ge\mu(k)$ implies $|s_m-s_n|\le\eps_k$. A regular sequence satisfies $|s_m-s_n|\le\eps_m+\eps_n$ uniformly for \emph{all} indices; then $\mu(k)=k+1$ is a modulus because $m,n\ge k+1$ yields $\eps_m+\eps_n\le2\eps_{k+1}=\eps_k$. Conversely a modulus yields the bare Cauchy property by choosing $k$ with $\eps_k<\delta$. Extracting a modulus function from a family of existential tails is a separate choice problem; it is not part of the bare Cauchy definition.

\subsection{Resampling and quotient equivalence}\label{supp:12-3}
\par\noindent{\small\emph{Related main-text discussion:} \hyperref[main:sub-41]{Section~\ref*{main:sub-41}: Regularization and equivalence}.}\par\smallskip

Set $\rho_\mu(n)=\max_{i\le n}\mu(i+1)$ and $r_n=s_{\rho_\mu(n)}$. If $m\le n$, both selected indices are beyond $\mu(m+1)$, so $|r_m-r_n|\le\eps_{m+1}\le\eps_m+\eps_n$; the other order of $m,n$ is symmetric. For $n\ge\max(k,\mu(k+1))$, the original index $n$ and $\rho_\mu(n)$ lie beyond $\mu(k+1)$, showing $|r_n-s_n|\le\eps_{k+1}\le\eps_k$. Hence resampling produces a regular equivalent sequence. The equivalence relation is transitive by intersecting two tails at precision $k+1$ and adding their errors. The quotient $\RR(K)$ is defined from \emph{modulated} data, not from all sequences satisfying only the bare Cauchy predicate.

\subsection{Ring congruence and the rational embedding}\label{supp:12-4}
\par\noindent{\small\emph{Related main-text discussion:} \hyperref[main:sub-42]{Section~\ref*{main:sub-42}: Addition, negation and products}.}\par\smallskip

Sum congruence follows by decomposing $(s+t)-(s'+t')=(s-s')+(t-t')$ and taking equivalence tails at precision $k+1$. Negation preserves absolute differences. For products, obtain eventual rational bounds $|s_n|\le 2^{B_s}$ and $|t'_n|\le2^{B_{t'}}$ \emph{independently}, using bounded tail values and the respective Cauchy estimates. In particular the bound for $t$ cannot automatically be reused for its equivalent representative $t'$. For indices beyond both bounding tails and the equivalence tails,
\[
|s_nt_n-s'_nt'_n|\le |s_n|\,|t_n-t'_n|+|t'_n|\,|s_n-s'_n|.
\]
Choose the $t\sim t'$-tail at precision $k+1+B_s$ and the $s\sim s'$-tail at precision $k+1+B_{t'}$. The two terms then do not exceed $2^{B_s}\eps_{k+1+B_s}=\eps_{k+1}$ and $2^{B_{t'}}\eps_{k+1+B_{t'}}=\eps_{k+1}$, respectively; their sum is $\eps_k$. Applied separately to two Cauchy sequences with their own eventual bounds, the same estimate constructs the product modulus. Hence all operations descend to the quotient. The ring laws follow pointwise, and constant sequences embed the rational base as a ring. If constant sequences $a,b$ are equivalent, $|a-b|\le\eps_k$ for all $k$; choosing $\eps_k<|a-b|$ when $a\ne b$ gives a contradiction, so the embedding is injective.

\subsection{Nonnegative classes, strict order and apartness}\label{supp:12-5}
\par\noindent{\small\emph{Related main-text discussion:} \hyperref[main:sub-43]{Section~\ref*{main:sub-43}: Positive cone, weak order and apartness}.}\par\smallskip

Sequence positivity means $\exists k,N\ \forall n\ge N: s_n\ge\eps_k$, while sequence nonnegativity means $\forall k\ \exists N\ \forall n\ge N:s_n\ge-\eps_k$. At quotient level, a \emph{single} representative satisfying the entire predicate is supplied. Equivalent representatives preserve positivity by halving a strict gap and preserve approximate nonnegativity by allocating half the error budget to each sequence. Set $x<y$ iff $y-x$ is positive, $x\le y$ iff $y-x$ is nonnegative, and $x\Apt y$ iff $x<y$ or $y<x$. Nonnegative classes are closed under addition using errors $\eps_{k+1}$; strict positives add by retaining a positive fixed lower bound. From $x\le y$ and $y\le x$, the difference is eventually between $-\eps_k$ and $\eps_k$ for every $k$, hence equivalent to zero, yielding antisymmetry. For strict cotransitivity $x<y$, choose a rational index where $y-x$ has positive lower bound $\eps_k$ and decide whether the approximation to a third element lies above or below the rational midpoint gap. Errors $\eps_{k+3}$ yield respectively $x<z$ or $z<y$. Tightness of apartness is derived from the constructive negative-of-strict-positivity characterization; trichotomy is \emph{not} asserted.

\subsection{Apartness-based inversion and ordinary inequality}\label{supp:12-6}
\par\noindent{\small\emph{Related main-text discussion:} \hyperref[main:sub-44]{Section~\ref*{main:sub-44}: Inverses for apart elements}.}\par\smallskip
\label{expanded:inverse}
A positive separation witness for a representative $s$ consists of two indices $k,N$ with $s_n\ge\eps_k>0$ for every $n\ge N$. Instead of assigning arbitrary values before $N$, the actual construction \emph{resamples} into the positive tail:
\[
 t_n=\bigl(s_{N+n}\bigr)^{-1}.
\]
All denominators are strictly positive. Their reciprocals satisfy $0<t_n\le2^k$ by the order-reversing inverse inequality. For nonzero rationals $a,b$ one has
\[
 a^{-1}-b^{-1}=(b-a)a^{-1}b^{-1}.
\]
Hence, whenever $m,n\ge\mu_s(j+2k)$, the Cauchy modulus for $s$ and the bound for both reciprocal factors give
\[
 |t_m-t_n|\le \eps_{j+2k}2^k2^k=\eps_j.
\]
Thus $\mu_t(j)=\mu_s(j+2k)$ is a concrete modulus (the source writes $j+(k+k)$). The quotient product with the \emph{unshifted} $s$ still equals one: the difference
\[
 s_nt_n-1=\bigl(s_n-s_{N+n}\bigr)\bigl(s_{N+n}\bigr)^{-1}
\]
is eventually at most $\eps_{j+k}2^k=\eps_j$ in absolute value. This is the nontrivial part of proving $[s][t]=1$, since the two factors use different sample indices. A negative separation witness is treated by applying the positive case to $-s$ and negating the constructed reciprocal. The inverse is independent of the witness by cancellation against $1$. The resulting existence assertion requires $x\Apt0$, not only $x\ne0$. The project proves multiplication cancellation from mere non-equality, but producing an inverse under that weaker hypothesis is equivalent to an additional Markov principle.

\subsection{Quantitative density, bounds and arity transport}\label{supp:12-7}
\par\noindent{\small\emph{Related main-text discussion:} \hyperref[main:sub-45]{Section~\ref*{main:sub-45}: Density, Archimedean control and arity transfer}.}\par\smallskip

For $x<y$, choose representatives $a,b$ and a strict positive gap $b_n-a_n\ge\eps_k$ beyond some index. Choose an index $N$ also beyond their tails at precision $k+2$, and set $q=a_N+\eps_{k+1}$. The error estimate $2\eps_{k+2}=\eps_{k+1}$ ensures that the constant rational class $\iota(q)$ stays strictly between $x$ and $y$ with a remaining positive precision margin. An Archimedean bound for $\RR(K)$ is obtained by bounding one rational representative on a tail and enlarging the bound to dominate its finite initial section. An isomorphism of admissible rational coefficient fields, acting termwise on Cauchy data, preserves Cauchy moduli and equivalence because it preserves the rational order and absolute values; it therefore induces the reported arity comparisons of completions. These are specific transfers, not a universal identification with a classical real-number field.

\subsection{Convergence and uniqueness of limits}\label{supp:13-1}
\par\noindent{\small\emph{Related main-text discussion:} \hyperref[main:sub-46]{Section~\ref*{main:sub-46}: Convergence and its uniqueness}.}\par\smallskip

Let $c=(s,\mu)$ be modulated rational Cauchy data and $x=[c]$. The constant rational class $\iota(s_n)$ approaches $x$: if $n\ge\mu(k+1)$, then for all later $m$, $|s_n-s_m|\le\eps_{k+1}$, and passing to the quotient gives the prescribed $\eps_k$-neighborhood convergence estimate. If a real sequence converges to both $x$ and $y$, at any precision $k$ choose a common tail on which its terms lie within $\eps_{k+1}$ of both limits. The triangle inequality bounds the separation of $x,y$ by $\eps_k$ for every $k$, and the established tightness/antisymmetry rules give $x=y$. This is the uniqueness theorem for the project's stated convergence predicate.

\subsection{The tracked diagonal limit}\label{supp:13-2}
\par\noindent{\small\emph{Related main-text discussion:} \hyperref[main:sub-47]{Section~\ref*{main:sub-47}: Tracked diagonal limits}.}\par\smallskip

The actual rational tracking datum contains \emph{one} rational sequence $q:\N\to K$ and one precision schedule $M$, not a two-index family of rational approximants. Its defining condition is
\[
i\ge M(k)\quad\Longrightarrow\quad
\mathsf{RClose}\bigl(f_i,\iota(q_i),\eps_k\bigr).
\]
Suppose also that the real sequence $f$ has the supplied Cauchy modulus $\nu$. Define
\[
\lambda(k)=\max\{\nu(k+2),M(k+2)\}.
\]
For $m,n\ge\lambda(k)$, the two approximation inequalities around $f_m,f_n$ and the real Cauchy estimate concatenate into
\[
\mathsf{RClose}\bigl(\iota(q_m),\iota(q_n),
\eps_{k+2}+\eps_{k+2}+\eps_{k+2}\bigr).
\]
Reflection of rational closeness along the injective ordered embedding then gives
\[
|q_m-q_n|\le3\eps_{k+2}\le\eps_k.
\]
Consequently the \emph{given sequence $q$ itself}, equipped with the explicit modulus $\lambda$, is modulated rational Cauchy data; no new diagonal choice of the $q_i$ is needed. Let $L=[q]$. To prove convergence of $f$, use the index threshold
\[
N(k)=\max\{\lambda(k+1),M(k+1)\}.
\]
For $i\ge N(k)$, tracking gives $\mathsf{RClose}(f_i,\iota(q_i),\eps_{k+1})$ and convergence of embedded rational terms gives $\mathsf{RClose}(\iota(q_i),L,\eps_{k+1})$. Transitivity produces $\mathsf{RClose}(f_i,L,\eps_k)$. This is precisely the source's $\mathsf{approxMod}$, $\mathsf{approxSeq}$ and $\mathsf{limit\_of\_approx}$ construction, with no countable representative selection hidden in the proof.

\subsection{Finite algebraic expressions of tracked families}\label{supp:13-3}
\par\noindent{\small\emph{Related main-text discussion:} \hyperref[main:sub-47]{Section~\ref*{main:sub-47}: Tracked diagonal limits}.}\par\smallskip

A tracked sequence $x_n$ comes with a specified limit and rational approximation data. For two tracked sequences $x_n\to x$, $y_n\to y$, the difference $(x_n+y_n)-(x+y)$ splits into the sum of two errors, so choosing each at $\eps_{k+1}$ gives the $\eps_k$ estimate. Negation preserves errors. For products, the same two-term identity used for rational Cauchy products applies to approximated real terms, after obtaining eventual absolute bounds for both sequences from their tracking data. This construction supplies closure under each of the finite operations. Structural induction on the finite syntax of expressions generated by constants, variables, negation, addition and multiplication then shows that evaluation commutes with the tracked limit. It does not justify arbitrary infinite expressions or unrestricted functional calculus.

\subsection{Modulated completeness versus rational approximation}\label{supp:13-4}
\par\noindent{\small\emph{Related main-text discussion:} \hyperref[main:sub-48]{Section~\ref*{main:sub-48}: Exact characterization of modulated completeness}.}\par\smallskip

The formal principle $\mathsf{UnrestrictedCompleteness}(K)$ quantifies over real Cauchy sequences \emph{already equipped with a modulus}. The principle $\mathsf{CauchyApprox}(K)$ supplies the rational approximation schedule needed for diagonal extraction. The source proves their equivalence. In the direction $\mathsf{CauchyApprox}\Rightarrow\mathsf{UnrestrictedCompleteness}$, the supplied rational approximants and real modulus form tracked data and the diagonal theorem gives a limit. In the reverse direction, let $(f,\nu)$ be the supplied modulated real Cauchy family. Apply completeness to obtain a limit $L$. Choose \emph{one} quotient representative $c$ of this single limit (quotient elimination, not countable representative selection), and set $q_i=\operatorname{approxAt}(c,i)$ and
\[
M(k)=\max\{\nu(k+1),k+1\}.
\]
For $i\ge M(k)$, the Cauchy-and-convergence estimate gives $f_i$ within $\eps_{k+1}$ of $L$, while the rational approximation lemma gives $L$ within $\eps_i\le\eps_{k+1}$ of $\iota(q_i)$. Combining the two closeness estimates and using $2\eps_{k+1}=\eps_k$ yields the required $\mathsf{ApproxData}(f,q,M)$. No countable choice of quotient representatives is used in this reversal. The manuscript does not interpret ``unrestricted'' as modulus-free, nor does it claim that a mere existence of a suitable rational approximant at each index automatically supplies one choice function.

\subsection{What choice supplies, and what it does not}\label{supp:13-5}
\par\noindent{\small\emph{Related main-text discussion:} \hyperref[main:sub-49]{Section~\ref*{main:sub-49}: Choice and modulus-free statements}.}\par\smallskip

The proposition $\mathsf{ChoiceNN}$ asserts countable selection for families of propositions over the unary index set. For each real class $x_i$, dyadic density gives a pair of path-count numerators $a_i,b_i$ with
\[
\mathsf{RClose}\bigl(x_i,\iota((a_i-b_i)\eps_{i+1}),\eps_i\bigr).
\]
The proof first produces such a pair only existentially for each $i$: it uses the Archimedean bound to bracket a rational approximant between successive unary integers after multiplying by $2^{i+1}$. The source applies $\mathsf{ChoiceNN}$ twice, first to select $a_i$ together with the existence of a suitable $b_i$, then to select $b_i$ for the previously chosen $a_i$. This yields a \emph{single} function $q_i=(a_i-b_i)\eps_{i+1}$ and the pointwise approximation theorem $\mathsf{SequenceApprox}(K)$. Together with the given real Cauchy modulus these data feed into the diagonal-limit construction. Under this explicit premise the source obtains modulated completeness. With additional modulus-selection machinery it also obtains the stated modulus-free consequence. These are \emph{conditional} implications with the choice proposition visible in their hypotheses; $\mathsf{ChoiceNN}$ has not become a globally assumed axiom. No converse to this choice implication and no metamathematical nonderivability claim is implied.

\subsection{Suprema from supplied rational upper data}\label{supp:13-6}
\par\noindent{\small\emph{Related main-text discussion:} \hyperref[main:sub-50]{Section~\ref*{main:sub-50}: Supplied data for least upper bounds}.}\par\smallskip

For an inhabited predicate $S\subseteq\RR(K)$, $\mathsf{SupData}(S)$ gives a rational upper bound $u_k$ and a point $s_k\in S$ with $\iota(u_k)\le s_k+\iota(\eps_k)$. Comparing $u_k$ with $s_m$ and vice versa gives $-\eps_k\le u_m-u_k\le\eps_m$, hence $|u_m-u_k|\le\eps_m+\eps_k$. Therefore $u$ is a regular rational sequence defining $U=[u]$. Each $u_k$ is an upper bound for all of $S$; passing to the quotient limit and using closure under arbitrarily small errors gives $x\le U$ for all $x\in S$. If $y$ is any upper bound, the near-attainment witnesses yield $\iota(u_k)\le y+\iota(\eps_k)$ for every $k$, so $U\le y$. Antisymmetry gives uniqueness of the least upper bound. The essential hypothesis is the \emph{joint family of rational bounds and nearly attaining elements}, not mere boundedness.

\subsection{Uniformly continuous extension with explicit modulus}\label{supp:13-7}
\par\noindent{\small\emph{Related main-text discussion:} \hyperref[main:sub-52]{Section~\ref*{main:sub-52}: Extension of uniformly continuous rational maps}.}\par\smallskip

For $f:K\to K$ and a function $\omega:\N\to\N$, the source's modulus states that $|x-y|\le\eps_{\omega(k)}$ entails $|f(x)-f(y)|\le\eps_k$. If $c=(s,\mu)$, then $f\circ s$ has modulus $k\mapsto\mu(\omega(k))$. The same implication maps equivalent rational sequences to equivalent image sequences, hence induces $\widehat f:\RR(K)\to\RR(K)$. For a constant rational sequence, $\widehat f(\iota(q))=\iota(f(q))$. If another map agrees on embedded rationals and preserves convergence along the rational approximants of \emph{every} representative (the precise $\mathsf{RatSeqContinuous}$ predicate), apply limit uniqueness to the two identical sequences of rational images to show the maps agree at each quotient class. Moduli compose as $\omega_f\circ\omega_g$ in the correct order, giving the extension-of-composition theorem. The target is $\RR(K)$; no general universal property for arbitrary complete codomains is proved.

\subsection{Markov's principle and inversion from inequality}\label{supp:14-1}
\par\noindent{\small\emph{Related main-text discussion:} \hyperref[main:sub-53]{Section~\ref*{main:sub-53}: Boolean-sequence principles}.}\par\smallskip

Encode a Boolean stream by the nonnegative modulated rational sequence whose value is $\eps_j$ after detection of its first success at $j$, and zero beforehand. The modulus $\mu(k)=k$ is valid because any newly detected success after stage $k$ has value at most $\eps_k$. The resulting class equals zero precisely when the stream is everywhere false. If all rational gap tests $3\eps_n<|r_n|$ on a regular representative $r$ fail, then $|r_n|\le3\eps_n$, so the class is zero. From $x\ne0$ we infer \emph{double-negated} success of a gap test; the project's Boolean MP supplies an index witnessing success, hence apartness and the existing apartness-based inverse. Conversely, assume all nonzero classes possess inverses; such an inverse forces apartness by the source's multiplication-order result. For a double-negated successful stream, the encoding is nonzero, hence apart, and its positive gap yields a detected success. This proves exactly the stated equivalence, without proving MP itself.

\subsection{LPO, trichotomy and zero-or-apartness}\label{supp:14-2}
\par\noindent{\small\emph{Related main-text discussion:} \hyperref[main:sub-54]{Section~\ref*{main:sub-54}: From a Boolean stream to a rational Cauchy encoding}.}\par\smallskip

For any regular representative $r$, the rational gap test $G(n):3\eps_n<|r_n|$ is decidable using the base field's decision procedure. Applying LPO gives a successful index, which creates apartness by regularity, or failure at every index, which implies equality of its quotient class to zero. Thus LPO gives $x=0\lor x\Apt0$. Conversely apply this disjunction to the Boolean encoding of a stream: equality to zero forces all values false, whereas apartness forces some finite success; this is LPO. Applying the same principle to differences $y-x$ gives the source's trichotomy equivalence. These arguments use decidable \emph{rational} tests and an \emph{additional} LPO proposition, not a primitive decision procedure for all real comparisons.

\subsection{WLPO and decidable zero}\label{supp:14-3}
\par\noindent{\small\emph{Related main-text discussion:} \hyperref[main:sub-54]{Section~\ref*{main:sub-54}: From a Boolean stream to a rational Cauchy encoding}.}\par\smallskip

WLPO applied to the Boolean stream of rational gap tests decides whether \emph{all} tests fail, in which case the regular representative defines zero, or whether it is false that all tests fail, in which case the quotient class cannot equal zero. This gives decidability of equality to zero without producing a positive separation witness in the second branch. Conversely, for a Boolean stream $p$, decidability of equality to zero of its encoded Cauchy class decides whether all entries of $p$ are false or the negation of that statement. This is precisely WLPO, not LPO. The distinction between ``not all false'' and a supplied successful index is essential.

\subsection{One-way logical bounds for strong principles}\label{supp:14-4}
\par\noindent{\small\emph{Related main-text discussion:} \hyperref[main:sub-55]{Section~\ref*{main:sub-55}: Lower bounds for stronger completeness principles}.}\par\smallskip

Assume every inhabited bounded predicate of real classes has a least upper bound. For a stream $p$, consider the predicate containing $0$ and also $1$ if there exists any success. This predicate is inhabited and bounded. The two test points may be chosen as $\eps_2<\eps_1$. Strict cotransitivity at the alleged least upper bound $y$ gives $y<\eps_1$ or $\eps_2<y$. In the first branch, a successful stream would put $1$ in the predicate, hence $1\le y<\eps_1$, impossible; thus every entry is false. In the second branch, if every entry were false, then $S=\{0\}$ and $0$ would be an upper bound, forcing $y\le0$ by minimality, contrary to $\eps_2<y$. Thus the all-false statement is refuted. These branches provide exactly WLPO (their order may be exchanged). Similarly, for bounded monotone convergence use $a_n=1$ after the first success detected by stage $n$ and $a_n=0$ earlier. The comparisons can be made quantitatively. Let $\eps_1=\frac12$, $\eps_2=\frac14$ and $\eps_3=\frac18$. Strict cotransitivity of $\eps_2<\eps_1$ at the candidate limit $L$ gives either $L<\eps_1$ or $\eps_2<L$. In the latter case use convergence at precision $\eps_3$ to choose a stage $N$ for which $L\le a_N+\eps_3$. If $a_N=0$, then $L\le\eps_3<\eps_2$, a contradiction; therefore $a_N=1$ and a finite success has occurred. In the former case, if a success ever occurred then all subsequent $a_n=1$; convergence at precision $\eps_3$ would give $1\le L+\eps_3<\eps_1+\eps_3<1$, a contradiction. Therefore all values of the stream are false. Hence arbitrary bounded monotone convergence implies LPO. Only these implications are proved; converses and independence results are not supplied.

\subsection{Exact meaning of the constructive scope}\label{supp:14-5}
\par\noindent{\small\emph{Related main-text discussion:} \hyperref[main:sub-55]{Section~\ref*{main:sub-55}: Lower bounds for stronger completeness principles}.}\par\smallskip

The absence of a call to \texttt{Classical.choice} in a verified theorem's transitive axiom closure is an audit fact, distinct from a metamathematical theorem that some proposition cannot follow from the underlying logic. Built-in quotient soundness and propositional extensionality are explicitly admitted by the recorded policy. Principles such as MP, LPO, WLPO and $\mathsf{ChoiceNN}$ occur as \emph{premises} of conditional statements or as terms within equivalences, not as silently imported axioms. Consequently one must distinguish ``no unconditional instance has been proved in this library'' from ``mathematically impossible'' or ``independent of the background theory.''

\subsection{Identical rational values and varying ranks}\label{supp:15-1}
\par\noindent{\small\emph{Related main-text discussion:} \hyperref[main:sub-56]{Section~\ref*{main:sub-56}: Equal fraction fields and unequal coordinate ranks}.}\par\smallskip

For finite alphabets with at least two child letters, transport through the common auxiliary fraction field gives isomorphic represented rational fields. For positive arities $b,c$, however, $F_b$ and $F_c$ are distinct free modules with bases indexed by $\Fin(b)$ and $\Fin(c)$, and their transverse kernels have constructed bases with $b-1$ and $c-1$ members. Equality or isomorphism of the scalar fields does not identify these modules, so there is no conflict between arity-independent scalar arithmetic and arity-dependent local algebraic ranks. Nothing here computes a topological dimension.

\subsection{Why a path-count sequence cannot distinguish rays}\label{supp:15-2}
\par\noindent{\small\emph{Related main-text discussion:} \hyperref[main:sub-57]{Section~\ref*{main:sub-57}: Primitive observations do not classify rays}.}\par\smallskip

For any rays $\xi,\eta$ and any index $n$, $S(E(\xi|_n))=n=S(E(\eta|_n))$. Therefore the entire sequence of raw prefix counts is identical for all rays. In particular no injective labeling of the rays can be obtained by retaining only this count sequence, whether or not a limit could be assigned; moreover that count sequence is not itself modulated Cauchy in the rational field. This is a statement about a specific forgotten-data map. It does not forbid inserting other letter-dependent numerical weights or separately defining maps of the boundary.

\subsection{Three distinct limiting constructions}\label{supp:15-3}
\par\noindent{\small\emph{Related main-text discussion:} \hyperref[main:sub-58]{Section~\ref*{main:sub-58}: Three incompatible kinds of completion or limit}.}\par\smallskip

The algebraic directed union of transverse finite-support vectors is constructed by \emph{zero-padding linear embeddings}; the ray boundary is defined by \emph{coherent infinite choices of child symbols}; and the numerical Cauchy quotient identifies \emph{modulated rational sequences} under asymptotic equality. Their elements, equality predicates and operations are distinct by construction. A map between any two of them would have to be defined and verified, including any claimed preservation of topology, algebra or convergence. No such blanket canonical identification or universal impossibility theorem is among the results recorded here.

\subsection{Root and endpoint level markers}\label{supp:15-4}
\par\noindent{\small\emph{Related main-text discussion:} \hyperref[main:sub-59]{Section~\ref*{main:sub-59}: Formal level labels}.}\par\smallskip

Word depth counts the number of edges on the finite root path, so the root is at depth zero. The source's extended-level function assigns the root the finite level marker $\mathrm{finite}(0)$; it assigns a ray the conventional symbol $+\infty$. The earlier suggested value $-\infty$ for the root is not the implemented convention. These assignments are labels in a particular extended-level type, not metrics, lengths of infinite paths or a theorem identifying rays with numerical infinities.

\subsection{Universal properties versus exhibited maps}\label{supp:16-1}
\par\noindent{\small\emph{Related main-text discussion:} \hyperref[main:sub-60]{Section~\ref*{main:sub-60}: Universal algebraic characterizations}.}\par\smallskip

A concrete isomorphism between the path-count semiring and a unary carrier does not by itself prove an initial-object property among all semirings. Likewise a normalized integer representation, a fraction-field comparison and an extension of uniformly continuous rational self-maps have their explicit universal quantifiers and target categories. General Grothendieck, localization or completion universal properties have not been established as project theorems. When one uses familiar universal-property language informally, it must be labeled external mathematical background rather than a claimed consequence of a particular checked declaration.

\subsection{Moduli, representatives and choice}\label{supp:16-2}
\par\noindent{\small\emph{Related main-text discussion:} \hyperref[main:sub-61]{Section~\ref*{main:sub-61}: Constructive completeness and choice}.}\par\smallskip

The existence of a Cauchy tail separately for every rational precision gives $\forall k\,\exists N$; a supplied modulus requires a single function $\mu$ with $\forall k$-indexed tail bounds. Constructively these are not interchangeable without a selection principle. At the quotient level a sequence of real classes similarly does not automatically provide one rational representative or a rational approximation array for \emph{each} index. The formal equivalence between modulated completeness and approximation data identifies an exact missing input; it is not a proof of impossibility or independence without that input.

\subsection{Classical field and order properties}\label{supp:16-3}
\par\noindent{\small\emph{Related main-text discussion:} \hyperref[main:sub-62]{Section~\ref*{main:sub-62}: Order, inverses and suprema}.}\par\smallskip

A constructive commutative ring equipped with a tight apartness relation may admit inverses for apart elements while failing to supply an inverse from mere negated equality. Universal inversion from $x\ne0$ is equivalent in this development to its Boolean Markov principle. Zero-or-apartness and trichotomy are tied to LPO, decidability of zero to WLPO, and unrestricted least upper bounds imply WLPO. None of these logical principles is added globally. Therefore one may not call $\RR(K)$ an unconditionally classically complete ordered field or assert a Dedekind-supremum theorem for every inhabited bounded subset.

\subsection{Scope of the longitudinal obstruction}\label{supp:16-4}
\par\noindent{\small\emph{Related main-text discussion:} \hyperref[main:sub-63]{Section~\ref*{main:sub-63}: Unbounded branching and topology}.}\par\smallskip

The proved contradiction uses the new coordinate of a zero-padded uniform finite vector together with an old coordinate. Its proof is insensitive to the topology because the ambient countable carrier consists solely of finitely supported sequences. A meaningful alternative longitudinal limit would require first specifying a larger ambient space, embeddings compatible with the chosen uniform directions, and a notion of convergence or algebraic colimit. The original obstruction does not refute every such unformulated alternative construction.

\subsection{Boundary compactness and charts}\label{supp:16-5}
\par\noindent{\small\emph{Related main-text discussion:} \hyperref[main:sub-63]{Section~\ref*{main:sub-63}: Unbounded branching and topology}.}\par\smallskip

For the countable alphabet the next-letter cylinder cover defeats finite subcovers and hence local compactness. Finite-alphabet boundary compactness may be a familiar classical fact under additional metatheoretic assumptions, but the source does not contain a general constructively formulated finite-cover theorem for it. Likewise the realline-chart obstruction assumes connectedness of appropriate target intervals; its statement cannot be strengthened to arbitrary positive-dimensional manifolds in the absence of corresponding chart models and target connectedness results.

\subsection{Remaining automorphism classification problem}\label{supp:16-6}
\par\noindent{\small\emph{Related main-text discussion:} \hyperref[main:sub-63]{Section~\ref*{main:sub-63}: Unbounded branching and topology}.}\par\smallskip

The levelwise permutation family is explicitly constructed and acts transitively on rays, while a single subtree-supported involution proves that the family does not exhaust all root-preserving automorphisms. This supplies both a positive subgroup result and a negative \emph{specific} completeness counterexample. It leaves open, in this development, a full structural description of the group of all root-preserving incidence automorphisms; a classification cannot be inferred merely from the existence of one additional element.

\subsection{Distinguishing theorem status precisely}\label{supp:16-7}
\par\noindent{\small\emph{Related main-text discussion:} \hyperref[sec:limitations]{Section~\ref*{sec:limitations}: Limitations and open mathematical questions} (contextual discussion; no separate main-text subsection).}\par\smallskip

There are five different cases that must remain distinct: (i) a theorem proved under no extra proposition beyond the displayed structural typeclasses, (ii) a theorem assuming an additional mathematical proposition such as $\mathsf{ChoiceNN}$, (iii) a formally constructed counterexample to a specific statement, (iv) a well-defined but not proved stronger proposition, and (v) an interpretation for which even the mathematical object or map has not been specified. The fifth is not an instance of the fourth, and an absent Lean proof alone never establishes metamathematical independence.

\subsection{Established results at their actual hypotheses}\label{supp:17-1}
\par\noindent{\small\emph{Related main-text discussion:} \hyperref[sec:conclusion]{Section~\ref*{sec:conclusion}: Conclusion} (contextual discussion; no separate main-text subsection).}\par\smallskip

The coherent collection consists of: unary word arithmetic and its equal-depth quotients under inhabitedness; normalized integer differences; triple representatives of the internally defined rational fraction field when two distinct letters are supplied; an explicit $1+(b-1)$ local coordinate decomposition over a characteristic-zero coefficient field; complementary bilinear projections; a countably indexed algebraic transverse basis and the zero-padding obstruction for uniform lines; native ray-boundary cylinder topology and automorphism actions; and a constructive rational Cauchy quotient with apartness inverses, tracked limits, and conditional stronger principles. Each belongs to the carrier and hypotheses where it was proved.

\subsection{Open directions without stronger claims}\label{supp:17-2}
\par\noindent{\small\emph{Related main-text discussion:} \hyperref[sec:conclusion]{Section~\ref*{sec:conclusion}: Conclusion} (contextual discussion; no separate main-text subsection).}\par\smallskip

General completeness without moduli or chosen rational approximation arrays, ordinary inversion from non-equality without an additional logical principle, general suprema, possible alternative longitudinal limit objects, complete automorphism classifications, and canonical boundary--number comparisons remain outside the unconditional theorem list. Their investigation is mathematically legitimate, but neither a prior theorem nor a documented proof gap should be recast as a no-go theorem about every possible construction. The correct conclusion is the explicitly delimited body of mathematics developed in this paper.
